\documentclass[11pt]{article}
\usepackage{amsmath,amssymb,amsthm}
\newtheorem{proposition}{Proposition}
\title{Exact-minimum alignment: a sharper connected response under original block composition}\author{}\date{}
\begin{document}\maketitle
\paragraph{Input class.}
Use the shared face posterior and common background calculation: the actual shared-face posterior, all-integer radial recursion,
common-forest harmonic naturality and retained $R^{15}$ compact-tail cost.
We work in its COMMON relative forest, with actual nonflat boundary $b_t=b+tv$,
$\|v\|\le1$, $\alpha=\sqrt{N/B}$. The admitted block data are
\[
 \|b_t\|\le CR^8,\quad \|m_t\|\le CR^8,\quad E_{\rm quad}(b_t)\le CR^{10},
 \quad K_{JJ}\ge cR^{-6},\quad n_J\le CR^4.                              \tag{1}
\]
The original outer prepared blocks of the boundary only preparation and nonflat block calculation are included. Every compact
interior is integrated without a witness selector. The chart radius is
$cR^{-6}$ in joint angular norm. The original normalized complement satisfies
\[
 q\le\exp[-cBR^{-18}+CR^{15}+CR^4\log(2+BR)].                            \tag{2}
\]
All boundary-only preparation frames are fixed from the starting record and
are not recomputed from the integration variables.

\paragraph{1. The actual minimum exists uniquely in this admitted class.}
The original ACTION, not a reference density, is strongly convex in the local
canonical chart with floor $cR^{-6}$. At its quadratic harmonic center $m_t$
its gradient is at most $CR^{16}/\sqrt B$. Consequently its local critical
point $x_t^*$ obeys
\[
 \|x_t^*-m_t\|\le CR^{22}/\sqrt B.                                     \tag{3}
\]
For $B\gg R^{34}$ this is at most $CR^5$, much smaller than the available
chart radius. Existence follows, for example, by strong convexity on a ball
about $m_t$ whose outward radial derivative is positive; uniqueness follows
from the same strict convexity.
The action at $m_t$ is at most $CR^{10}$, including its actual cubic remainder
$CR^{24}/\sqrt B$. Outside the chart the geometric action bound is
$cBR^{-18}-CR^{15}$. For $B\gg R^{34}$ this is larger than that reference
value. Hence the critical point is the UNIQUE GLOBAL compact minimum.
It depends smoothly on $t$ by the nonsingular interior Hessian. Write its
actual group links as $U^*(t)$ and its exact classical action as
\[
 V(t)=S_B(U^*(t);\eta(t)).                                               \tag{4}
\]
This is not a Gaussian vacuum approximation.

The actual original law has centered moments around this action minimum:
\[
 E\|z\|^{2p}\le C_pR^{10p}+\text{compact tail},\qquad
 U_j=e^{\alpha z_j}U_j^*(t).                                            \tag{5}
\]
Indeed the shared face posterior and common background calculation gives the same moment bound around $m_t$, (3) is at most $CR^5$,
and the two near-identity group coordinate systems have uniformly comparable
distances. The true Haar density remains in the law; its minimum need not
coincide with the action minimum.

\paragraph{2. An explicit Haar flow following that minimum.}
Set
\[
 \Phi_t(U)_j=U_j^*(t)U_j^*(0)^{-1}U_j,\qquad
 \dot U_j^*(t)U_j^*(t)^{-1}=\alpha w_j^*(t).                             \tag{6}
\]
At each fixed time this is a product of literal left Haar translations,
with Jacobian one and divergence zero. The reconstructed fine-link and
observable maps are their actual compositions with the common forest; the
boundary output is the declared $\eta(t)=\exp(\alpha b_t)$.
The canonical left boundary velocity is
$q_b(t)=\operatorname{dexp}_{\alpha b_t}v$; its norm is bounded and
$\|\dot q_b\|\le C\alpha$ for unit $v$.

Differentiating the interior stationarity equations at $U^*$ gives
\[
 H^*_{JJ}w^*+H^*_{Jb}q_b=0.                                             \tag{7}
\]
Here $H^*$ is the ACTUAL canonical left Hessian at the minimum. Interior
commutator ambiguities vanish because its interior gradient is zero.
The full Hessian differs from the flat complete incidence by operator norm
at most $C\alpha R^8$. Since $B\gg R^{28}$, this is smaller than half the
interior floor $cR^{-6}$. Also $\|H^*\|\le C$ and the full form is bounded
below by $-C\alpha R^8 I$. Apply this lower bound to $(q_b,w^*)$ and use(7):
\[
 \|w^*\|\le CR^3.                                                      \tag{8}
\]
Explicitly, $\kappa\|w^*\|^2\le w^{*T}H^*_{JJ}w^*
\le C+C\alpha R^8\|w^*\|^2$, so the last term is absorbed. Thus the actual
minimum-following flow has the same norm scale as the quadratic extension,
without assuming their velocities identical.

\paragraph{3. Exact removal of the linear fluctuation score.}
Let $D=\partial_t+W_t$, with $W_t=\alpha\sum w_j^*(t)X_j$, and $G_t=DS_B$.
The flow divergence is zero. The envelope theorem and differentiated
stationarity give the EXACT identities
\[
 G_t(U^*)=V'(t),\qquad \nabla_{\rm interior}G_t(U^*)=0.                  \tag{9}
\]
The second statement follows directly from(7), not from bounding separate
plaquette Hessians or discarding covariance terms.

On the common chart all third canonical derivatives of the original local
plaquette action have operator norm $CB^{-1/2}$, with finite incidence and
$N$ fixed. Differentiating the actual exponential-product coordinates adds
terms of the same size; their interior-force contribution vanishes at the
minimum. With (8), Taylor's integral formula around $z=0$ gives
\[
 |G_t(U)-V'(t)|\le C B^{-1/2}R^3\|z\|^2.                              \tag{10}
\]
This is the improvement over estimating $G$ by the absolute large background
norm $\|x\|^2\sim R^{16}$. Equations(5),(10) imply
\[
 \operatorname{Var}_tG_t\le CR^{26}/B+\mathcal E,\qquad
 |E_tG_t-V'(t)|\le CR^{13}/\sqrt B+\mathcal E,                           \tag{11}
\]
where one may take $\mathcal E=CBR^{16}q$ after increasing constants.
Global original group derivative bounds control the compact complement;
it is not removed from the fiber. The normalized pullback affinity is at least
$1-\frac18[CR^{26}/B+\mathcal E]$.

\paragraph{4. A connected free-energy norm relative to the EXACT classical state.}
Let
\[
 \Gamma(t)=-\log[Z(\eta(t))/Z(\eta(0))],\qquad
 F(t)=\Gamma(t)-[V(t)-V(0)].                                            \tag{12}
\]
Both normalizers and both classical actions are exact. Define $h=(q_b,w^*)$.
In $DG_t$, the time derivative of the interior velocity contributes
$\sum\dot w_j^*F_j$, with $F_j=\alpha X_jS_B$. Its expectation is ZERO by
original interior Haar integration by parts; its value at $U^*$ is also zero.
It has not been omitted pointwise or bounded by a large inverse-mode factor.
The remaining terms are the same-field Hessian $H(U)[h,h]$ and the boundary
acceleration term $\dot q_b\cdot F_b(U)$. Their differences from their values
at $U^*$ are at most $CB^{-1/2}R^6\|z\|$ and
$CB^{-1/2}\|z\|$, respectively, on the chart.
Using the exact identity $\Gamma''=E DG-\operatorname{Var}G$ and
$V''=DG(U^*)$, we obtain the sharper CONNECTED bound
\[
 |F''(t)|\le C[R^{11}/\sqrt B+R^{26}/B]+\mathcal E.                     \tag{13}
\]
The fluctuation current $F'(0)$ need not vanish at a nonflat record. It is
retained, with $|F'(t)|\le CR^{13}/\sqrt B+\mathcal E$ from(11).
Thus the four-term/centered remainder, or equivalently
$F(t)-tF'(0)$, has bound one half of(13) times $t^2$; it is not legitimate to
erase the current baseline when reporting the raw action.
The nonlinear classical Hessian $V''$ contains its exact Schur term and the
boundary acceleration/current partner. Comparing it to the flat Gaussian
form still costs $CR^{14}/\sqrt B$ in this admitted nonflat class.

\paragraph{5. Compatibility with the actual shared-face composition.}
If a union has a unique global minimum, fixing its shared face at that
minimum makes each conditional subblock minimum its literal restriction.
This is NONLINEAR minimum naturality, by the variational definition and
uniqueness, in the same common forest. The shared face admitted in the shared face posterior and common background calculation
lies within $CR^5$ of this actual minimum as well: use(3) and its original
posterior fluctuation bound. Along the connecting face path, (7)--(8) bound
the local minimum response by $CR^3$, preserving a local mean norm $CR^8$.
The joint classical action is stationary in the face at the union minimum;
its face Schur Hessian is bounded above by $C$ on this chart, with the small
coordinate-acceleration correction retained. Its energy increase is therefore
at most $C\|\delta F\|^2$, preserving the $CR^{10}$ energy class.

More generally the exact classical state composes by minimizing, and the
fluctuation state composes by integrating:
\[
 V_{\mathcal U}(\eta)=\min_F\{S_\Sigma(F;\eta)+V_1(F)+V_2(F)\},
\]
\[
 e^{-F_{\mathcal U}(\eta)}=
 \int e^{-[S_\Sigma+V_1+V_2-V_{\mathcal U}]}e^{-F_1(F)-F_2(F)}dF.         \tag{14}
\]
Here in(14) the absolute fluctuation state is $F_i=-\log Z_i-V_i$;
relative differences are formed only afterwards as in(12).
Equation(14) retains raw Haar factors, live classical/current terms and every
normalizer. It is not an independent-block Gaussian product. The earlier
exact covariance formula for the face integral remains necessary when
calculating the new current or Hessian.

\paragraph{6. New local range and the remaining physical obligation.}
The admission, uniqueness, centered-moment and tail requirements above only
need the conservative $B\gg R^{34}$. Thus, for example,
\[
 R\le B^{1/36}                                                        \tag{15}
\]
is valid for this exact-classical-relative class for large $B$.
The variance in(11) is then $O(B^{-5/18})$, the connected curvature in(13)
is $O(B^{-7/36})$, and the raw fluctuation-current/increment error is
$O(B^{-5/36})$. The compact tail is negligible. These claims do NOT extend
the old absolute Gaussian scalar-increment approximation to this range:
its live nonlinear classical baseline has deliberately been retained exactly.

The improvement from $R^{38}/B$ to $R^{26}/B$ is a genuine cancellation in an
ACTUAL common realization, not a new reference-law assumption. Nevertheless
(13) still grows under $R\mapsto LR$ at fixed $B$. It is not a strict
physical-scale contraction, and physical radius $aB^{1/36}$ still vanishes
when $B\asymp\log(1/a)$. No all-volume good-event, summable tiling error,
noncontractible holonomy replacement or fullS94 conclusion follows.
\end{document}
